% !PRINTED 463
\weberpage{463}
The three arguments of (2) are therefore the roots of equivalent quadratic forms.

We distinguish the three cases:

1. $D\equiv 0\ (\mathrm{mod}\ 4)$. Here $B$ is even, and $A$ and $C$ cannot both be even; of the three forms (3), (5) one therefore has two odd outer coefficients, while the other two have one even and one odd outer coefficient.

2. If $D\equiv 1\ (\mathrm{mod}\ 8)$, then $B$ is odd, $B^2\equiv 1\ (\mathrm{mod}\ 8)$, $AC\equiv 0\ (\mathrm{mod}\ 2)$, hence at least one of the two coefficients is even, and among the three forms (3), (5) there is one, but also only one, that has two even outer coefficients.

3. If $D\equiv 5\ (\mathrm{mod}\ 8)$, then $4AC\equiv 4\ (\mathrm{mod}\ 8)$ and the two outer coefficients $A$ and $C$ are \gesperrt{odd}.

If it now proves possible to set up an equation with integral coefficients which only one of the three roots (2) satisfies, it follows from this that this \gesperrt{one} root is rationally expressible by $j(\omega)$ and is therefore a class invariant.

Between the functions
\weq{6}{u = f(\omega)^{24}, \quad v = f\left(\frac{c+\partial\omega}{a}\right)^{24} \quad (c\equiv 0)\ (\mathrm{mod}\ 2)}
there exists, for an odd degree of transformation $n$ (\S\S\,73, 74), a transformation equation:
\[
\Phi_n(u,v) = 0,
\]
and the equation:
\weq{7}{\Phi_n(u,u) = 0}
is therefore satisfied, by \S\,53, only when for one of the quantities (6):
\weq{8}{\frac{c+\partial\omega}{a} = \frac{\gamma+\delta\omega}{\alpha+\beta\omega},}
in which
\weq{9}{\begin{aligned}
&\begin{pmatrix} \alpha, & \beta \\ \gamma, & \delta \end{pmatrix} \equiv \begin{pmatrix} 1, & 0 \\ 0, & 1 \end{pmatrix} \quad (\mathrm{mod}\ 2)\\
\text{or}\quad &\begin{pmatrix} \alpha, & \beta \\ \gamma, & \delta \end{pmatrix} \equiv \begin{pmatrix} 0, & 1 \\ -1, & 0 \end{pmatrix} \quad (\mathrm{mod}\ 2)
\end{aligned}}
is a linear transformation belonging to the first or to the second class (\S\,36).

The comparison of (8) with (3) leads, however, as above, to the conditions:
%CONT>>
